% TOP_PROBLEM: {"id":30007514,"problem_number":"LOCAL-30007514","title":"Expectations beyond rational expectations","category_id":21,"category":{"id":21,"name":"economics","display_name":"Economics","description":"Open economic theory statements, published research agendas, and proposed scoped specifications; question provenance and historical priority are qualified per record.","slug":"economics","order_index":21,"created_at":"2026-10-08T00:00:00Z"},"status":"research_question","external_url":null,"legacy_ids":[],"published":false,"statement":"Construct and validate a computationally tractable subjective-belief model that fits household forecasts and responds endogenously to policy changes.","economics":{"original_id":3,"field":"Business cycles and macro dynamics","kind":"design","formalization_basis":"adapted_research_question","source_status":"explicit_open","priority_status":"earliest_found","priority_note":"This is the earliest explicit relevant statement verified in this audit, not a claim of original priority; older literature and earlier versions may contain antecedents.","earliest_verified_reference":{"authors":["Benjamin Moll"],"title":"Research Agenda: Benjamin Moll","year":2020,"url":"https://economicdynamics.org/research-agenda-moll2020/","doi":null,"locator":"Section 4.1, “Heterogeneous-agent macro as a gateway to behavioral macro”","evidence_summary":"The author explicitly proposes empirically disciplined departures from rational expectations, with testable predictions rather than unrestricted behavioral assumptions."},"current_reference":{"authors":["Benjamin Moll"],"title":"The Trouble with Rational Expectations in Heterogeneous Agent Models: A Challenge for Macroeconomics","year":2025,"url":"https://arxiv.org/html/2508.20571v1","doi":"10.48550/arXiv.2508.20571","locator":"Introduction; Sections 3.2–3.4 and 4","evidence_summary":"The paper explicitly asks what should replace rational expectations and sets out computational tractability, empirical discipline, and endogenous responses of beliefs as desiderata. The retrieved HTML reports a 2026 revision; the arXiv identifier was first public in 2025."},"formalization_note":"The source gives the three qualitative desiderata. This finite-state architecture, tolerances, and validation target are an adapted design problem."},"statusQualification":"Published question or research agenda verified at the cited locator. The mathematical specification is adapted for this catalogue and includes additional assumptions; source publication does not establish first-ever priority or certify this exact model remains unsolved. The source gives the three qualitative desiderata. This finite-state architecture, tolerances, and validation target are an adapted design problem.","statusReviewedAt":"2026-10-08"}
% FIRST_POSTED: 2026-10-08
% ENTRY_KIND: statement_only
% !TeX program = lualatex
\documentclass[11pt]{article}
\usepackage{fontspec}
\usepackage[a4paper,margin=25mm]{geometry}
\usepackage{amsmath,amssymb,mathtools}
\usepackage{xurl}
\usepackage[hidelinks]{hyperref}
\newcommand{\sref}[2]{\hyperlink{source:#1:#2}{[S#2]}}
\setlength{\emergencystretch}{3em}
\begin{document}
% problemId:30007514
\section*{Expectations beyond rational expectations}\label{30007514}
\subsection{Definitions and mathematical statement}
Fix an incomplete-markets economy with heterogeneous households, idiosyncratic income risk, aggregate shocks, and a borrowing constraint. Let \(\mu_t\) denote the joint distribution of assets and income, \(s_t\) the aggregate fundamental state, and \(p_t\) the vector of equilibrium prices. Each household forecasts future prices through a subjective conditional distribution \(Q_\theta(\,\cdot\mid I_{it},b_{it})\), where \(I_{it}\) is its observed information and \(b_{it}\) is a finite-dimensional belief state. Specify a updating rule \(b_{i,t+1}=G_\theta(b_{it},I_{i,t+1},a_t)\), with \(a_t\) the observed policy regime. Expectations in household optimality conditions are taken under \(Q_\theta\); market clearing determines actual prices.
Construct and identify a pair \((Q_\theta,G_\theta)\) satisfying three jointly assessed requirements: computationally feasible global equilibrium with aggregate risk; agreement with held-out household forecast distributions and errors; and endogenous, testable changes in beliefs following policy changes. More precisely, fix tolerances \(\epsilon_F,\epsilon_M>0\) and require forecast-distance and macroeconomic prediction loss to fall below them on a declared evaluation sample. Demonstrate which restrictions identify \(\theta\) and prevent arbitrary fitting through unobserved beliefs. The target is a disciplined alternative to full rational expectations in this economy, with reproducible validation and sensitivity analysis, rather than proof that any one learning rule universally replaces rational expectations.

\par\textbf{Formulation and scope.} The source gives the three qualitative desiderata. This finite-state architecture, tolerances, and validation target are an adapted design problem.

\par\textbf{Open-question provenance.} An explicit question or research agenda was verified at the source locator. This does not by itself certify that every later version remains unresolved \sref{30007514}{1}.

\par\textbf{Historical priority.} This is the earliest explicit relevant statement verified in this audit, not a claim of original priority; older literature and earlier versions may contain antecedents.

\subsection{Short English statement}
Construct and validate a computationally tractable subjective-belief model that fits household forecasts and responds endogenously to policy changes.

\subsection{Sources}
\begin{itemize}\raggedright
\item[\textnormal{[S1]}]\hypertarget{source:30007514:1}{} Benjamin Moll. 2020. Research Agenda: Benjamin Moll. Section 4.1, “Heterogeneous-agent macro as a gateway to behavioral macro”.
\url{https://economicdynamics.org/research-agenda-moll2020/}.
{\small\textit{Earliest verified question/agenda reference; historical priority qualified below. The author explicitly proposes empirically disciplined departures from rational expectations, with testable predictions rather than unrestricted behavioral assumptions.}}

\item[\textnormal{[S2]}]\hypertarget{source:30007514:2}{} Benjamin Moll. 2025. The Trouble with Rational Expectations in Heterogeneous Agent Models: A Challenge for Macroeconomics. Introduction; Sections 3.2–3.4 and 4.
\url{https://arxiv.org/html/2508.20571v1}.
DOI: 10.48550/arXiv.2508.20571.
{\small\textit{Supporting or current literature; see evidence qualification. The paper explicitly asks what should replace rational expectations and sets out computational tractability, empirical discipline, and endogenous responses of beliefs as desiderata. The retrieved HTML reports a 2026 revision; the arXiv identifier was first public in 2025.}}
\end{itemize}
\end{document}
